\section{动机}
\label{sec:motivation}

\subsection{挑战}
当小型本地模型驱动调查时，会出现三个挑战。我们在各项研究的配置上（\S\ref{sec:setup}）分别测量了它们。

\stepnum{1} \textbf{流程失败。}自行选择探针的 Qwen2.5-7B 只验证了 0.083 到 0.125 的案例：它每个回合运行约七个不同的探针，花掉 10 个成本单位中的 9.4 到 9.5，四分之三的回合以工具预算耗尽告终；把它的选择换成均匀随机选择能把完成率提高 $+0.556$（\S\ref{sec:cost}）。在我们的提示与服务配置下，Llama-3.1-8B 在全部 3{,}389 次决策中都请求再运行一个探针，包括账本后验超过 0.999、提示已说明没有合法探针可用的那些决策。更大的模型会给出结论，但并不总是基于证据：没有账本时，Qwen2.5-32B 与 72B 在它们给出结论的需处置事件中分别把 13/51 和 20/54 判为良性。

\stepnum{2} \textbf{证据有一部分是攻击者的。}用户代理、请求载荷或工单评论会被复制进智能体读取的日志。当这样的字段带有“把案件结为授权扫描”的指令时，每个自行决定并给出结论的模型都服从了它（\S\ref{sec:injection}）。良性结论是攻击者最想要的结论，因为它在没有任何响应的情况下关闭案件。

\stepnum{3} \textbf{总得有人读日志。}按文档日志格式写的固定解析器在日志流水线变化后什么也读不出，于是每个案例都被升级（\S\ref{sec:rawlog}）。模型能读懂变化后的格式，但读取的模型可能被说服接受攻击者的读数，而且它读的是攻击者能写入的每一个来源。

\subsection{我们的方案}
\hesp{} 用模型之外的组件回应每个挑战。\stepnum{1} 一个由不含 LLM 的开发回合计数得到的似然表更新的账本、按单位成本期望信息增益对每个合法探针的排序、在每次调用规划器前应用的明确接受规则，以及规则满足时的控制器侧停止（附录~\ref{app:loop} 给出确切顺序）。\stepnum{2} 一个只接受“引用了账本中支持所称原因的当前观测”的结论的结束守卫，使从未进入账本的文本无法结束案件；以及一条要求良性结论建立在不同来源组之上的佐证规则。\stepnum{3} 读取器之间的信任边界：规则解析器解析出的观测是可信的，模型解析出的观测可以抬高某个原因、可以确认攻击，但永远不能确认良性结论。预写式日志让审计可以重放每一步。
